%  --- Définitions de chapitres ---
\newcolumntype{C}[1]{>{\centering\arraybackslash}m{#1}}
\newcolumntype{L}[1]{>{\raggedright\arraybackslash}m{#1}}
\newcolumntype{P}{C{0.05\linewidth} L{0.25\linewidth} L{0.25\linewidth} C{0.25\linewidth} L{0.11\linewidth}}

\newcommand{\phdr}{%
    \textbf{N\textsuperscript{o}} & \textbf{Points à mesurer} & \textbf{Valeur attendue} &
    \textbf{Valeur mesurée} & \textbf{OK/NOK} \\}

\newcounter{mesure}
\newcommand{\ydnm}{\stepcounter{mesure}\arabic{mesure}}
% ---------------------------------

\chapter{Protocole de test et de validation matérielle}
\label{ann:protocole-test}

\setlength{\tabcolsep}{4pt}
\renewcommand{\arraystretch}{2.5}

Ce protocole rassemble les mesures de mise en service de la
carte électronique réalisée. Les essais suivent une progression destinée à valider
le matériel par étapes successives : contrôles hors tension, alimentations,
communication et programmation, périphériques numériques, puis chaîne audio et
essais fonctionnels. Cet ordre répond à une exigence de sécurité : un étage
n'est sollicité qu'une fois son alimentation validée, afin de ne pas endommager
un composant sur un rail mal réglé ou court-circuité.

\paragraph{Précautions :}
La première mise sous tension s'effectue sur une alimentation de laboratoire à
courant limité, entre 200 et 300 mA, un
court-circuit échappé au contrôle à l'ohmmètre se traduit alors par une
limitation de courant plutôt que par une destruction. L'offset continu
en sortie de jack est vérifié avant le
branchement d'un casque, et les niveaux audio sont montés progressivement
depuis un volume numérique bas.

\paragraph{Matériel : }
Multimètre, alimentation de laboratoire à courant limité, oscilloscope, fichier
MP3 de test à 1~kHz, charge audio (casque de test ou résistances) et un analyseur
audio pour la mesure du THD+N et du SNR.

\paragraph{Lecture des tableaux : }
Pour les contrôles visuels, la valeur attendue désigne l'état conforme
recherché. La dernière colonne consigne le verdict (OK/NOK) et les remarques éventuelles.

\clearpage



% =====================================================================
\section{Contrôles avant mise sous tension}
% =====================================================================

Contrôles au ohmmètre, carte hors tension et \textbf{batterie déconnectée}, précautions ESD.
Pour chaque mesure rail–masse, laisser la lecture se stabiliser (charge des condensateurs de découplage)
avant de relever la valeur finale. Critère de court-circuit : R < 10 \Omega.

\begin{longtable}{P}
    \caption{Inspection visuelle et contrôles à l'ohmmètre (carte hors tension).}
    \label{tab:proto-preflight}                                                                                                \\
    \toprule \phdr \midrule \endfirsthead
    \multicolumn{5}{l}{\footnotesize\itshape Tableau \thetable{} — suite}                                                      \\
    \toprule \phdr \midrule \endhead
    \midrule \multicolumn{5}{r}{\footnotesize\itshape suite page suivante}                                                     \\ \endfoot
    \bottomrule \endlastfoot
    \ydnm & Inspection visuelle (soudures, ponts, composants) & Aucun pont de soudure, aucun composant manquant ni décalé &  & \\ \hline
    \ydnm & Polarité des composants                           & Repère pin 1 / cathode conforme                           &  & \\ \hline
    \ydnm & Connecteur USB-C                                  & Aucun pont entre broches adjacentes                       &  & \\ \hline
    \ydnm & TP GND (continuité entre points)                  & $R < 1~\Omega$ entre tous les points GND                  &  & \\ \hline
    \ydnm & TP VBUS $\rightarrow$ GND                         & $R > 5~\Omega$                                            &  & \\ \hline
    \ydnm & TP VSys $\rightarrow$ GND                         & $R > 5~\Omega$                                            &  & \\ \hline
    \ydnm & TP VBatP $\rightarrow$ GND                        & $R > 5~\Omega$                                            &  & \\ \hline
    \ydnm & TP +3.3V $\rightarrow$ GND                        & $R > 5~\Omega$                                            &  & \\ \hline
    \ydnm & AVDD / CPVDD / DVDD $\rightarrow$ GND (PCM5242)   & $R > 5~\Omega$ sur chaque pin                             &  & \\ \hline
    \ydnm & PVDD $\rightarrow$ GND (MAX97220)                 & $R > 5~\Omega$                                            &  & \\ \hline
    \ydnm & TP VBUS / VSys / +3.3V / VBatP (isolement mutuel) & $R > 5~\Omega$ entre chaque paire de rails                &  & \\ \hline
    \ydnm & Connecteur batterie                               & Pins +/$-$ conforme au schéma                             &  & \\ \hline
\end{longtable}


\clearpage
% =====================================================================
\section{Mise sous tension et alimentations}
% =====================================================================
\begin{longtable}{P}
    \caption{Alimentations et gestion d'énergie (sous tension, alimentation à courant limité).}
    \label{tab:proto-power}                                                                                                    \\
    \toprule \phdr \midrule \endfirsthead
    \multicolumn{5}{l}{\footnotesize\itshape Tableau \thetable{} — suite}                                                      \\
    \toprule \phdr \midrule \endhead
    \midrule \multicolumn{5}{r}{\footnotesize\itshape suite page suivante}                                                     \\ \endfoot
    \bottomrule \endlastfoot
    \ydnm & Courant d'appel à la mise sous tension (alim. limitée) & Pas de surintensité, consommation au repos faible    &  & \\ \hline
    \ydnm & TP VBUS                                                & 4,75--5,25~V                                         &  & \\ \hline
    \ydnm & TP VBatP                                               & 3,0--4,2~V selon l'état de charge                    &  & \\ \hline
    \ydnm & TP VSys, sans batterie, USB débranché                  & $\approx$ 0 V                                        &  & \\ \hline
    \ydnm & TP VSys, sans batterie, USB branché                    & $\approx$ VBUS                                       &  & \\ \hline
    \ydnm & TP VSys, avec batterie, USB débranché                  & $\approx$ VBAT                                       &  & \\ \hline
    \ydnm & TP VSys, avec batterie, USB branché                    & $\approx$ VBUS                                       &  & \\ \hline
    \ydnm & Courant de charge (sonde sur batterie)                 & Conforme à I\textsubscript{FAST} fixé par résistance &  & \\ \hline
    \ydnm & TP VBatP (fin de charge)                               & 4,20~V $\pm$50~mV                                    &  & \\ \hline
    \ydnm & TP +3.3V (sortie buck TPS62A01)                        & 3,30~V                                               &  & \\ \hline
    \ydnm & PCM5242 — AVDD (pin)                                   & 3,3~V                                                &  & \\ \hline
    \ydnm & PCM5242 — CPVDD (pin)                                  & 3,3~V                                                &  & \\ \hline
    \ydnm & PCM5242 — VNEG (pin)                                   & $-$3,3~V                                             &  & \\ \hline
    \ydnm & PCM5242 — DVDD (pin)                                   & 3,3~V                                                &  & \\ \hline
    \ydnm & PCM5242 — LDOO (pin, LDO interne)                      & 1,8~V                                                &  & \\ \hline
    \ydnm & MAX97220 — PVDD (pin)                                  & 3,3~V                                                &  & \\ \hline
    \ydnm & TP Supervisor                                          & $\approx$ VSys                                       &  & \\ \hline
\end{longtable}